%===============================================================================
% APPENDIX XX: METHOD-EVIDENCE
% Disciplinary Boundaries: What Evidence Can EBF Use?
%===============================================================================
% Category: METHOD
% Version: 1.0
% Date: January 2026
% Language: english
% Dependencies: LIT-LEARNING (XIX), LIT-CAMERER (T), CORE-WHERE (BBB)
%===============================================================================

\documentclass[11pt, a4paper]{article}
\usepackage[utf8]{inputenc}
\usepackage[english]{babel}
\usepackage[margin=1in]{geometry}
\usepackage{amsmath, amssymb, amsthm}
\usepackage{booktabs, array, multirow, tabularx}
\usepackage{xcolor}
\usepackage{tcolorbox}
\usepackage{hyperref}
\usepackage{enumitem}

\tcbuselibrary{breakable}

\title{Appendix XX: METHOD-EVIDENCE\\
Disciplinary Boundaries: Behavioral Economics, Behavioral Science, and Neuroeconomics}
\author{BEATRIX Research Group \\ FehrAdvice \& Partners AG}
\date{January 2026}

\begin{document}

\maketitle

%===============================================================================
% METADATA & OVERVIEW
%===============================================================================

\begin{tcolorbox}[colback=gray!10!white, colframe=gray!75!black, title=Appendix Metadata]
\begin{tabular}{@{}ll@{}}
\textbf{Code:} & XX (METHOD-EVIDENCE) \\
\textbf{Category:} & METHOD: Methodology \\
\textbf{Version:} & 1.0 (January 2026) \\
\textbf{Purpose:} & Define which evidence types EBF can use for which purposes \\
\textbf{Related Chapters:} & Ch. 4 (Methodology), Ch. 18 (Meta-Theory) \\
\textbf{Related COREs:} & BBB (CORE-WHERE), AU (CORE-AWARE) \\
\textbf{Related Appendices:} & XIX (LIT-LEARNING), T (LIT-CAMERER), AN (METHOD-LLMMC) \\
\end{tabular}
\end{tcolorbox}

%===============================================================================
% CROSS-REFERENCE MAP
%===============================================================================

\begin{tcolorbox}[colback=cyan!5!white, colframe=cyan!75!black, title=Cross-Reference Map]
\textbf{Incoming References (Who cites XX):}
\begin{itemize}[nosep]
    \item \textbf{PREDICT-HABIT (HF):} Evidence tier for $\theta_1$, $\theta_2$
    \item \textbf{LIT-LEARNING (XIX):} IV-identified habit parameters
    \item \textbf{CORE-WHERE (BBB):} Parameter source classification
    \item \textbf{METHOD-LLMMC (AN):} Evidence tier requirements for LLMMC input
    \item \textbf{All LIT Appendices:} Evidence quality assessment
\end{itemize}

\textbf{Outgoing References (What XX depends on):}
\begin{itemize}[nosep]
    \item \textbf{Appendix G (Glossary):} Definitions of methodological terms
    \item \textbf{Philosophy of Science:} Demarcation criteria
\end{itemize}
\end{tcolorbox}

%===============================================================================
% CHAPTER LINKAGE
%===============================================================================

\begin{tcolorbox}[colback=blue!5!white, colframe=blue!75!black, title=Chapter Linkage]
\textbf{Primary Chapter:} Chapter 4 (Methodology)
\begin{itemize}[nosep]
    \item Establishes methodological standards for EBF
    \item Defines evidence hierarchy
\end{itemize}

\textbf{Secondary Chapters:}
\begin{itemize}[nosep]
    \item Chapter 18 (Meta-Theory): Philosophical foundations
    \item Chapter 5 (Awareness): Uses learning evidence
\end{itemize}
\end{tcolorbox}

%===============================================================================
% QUICK REFERENCE
%===============================================================================

\begin{tcolorbox}[colback=yellow!5!white, colframe=yellow!75!black, title=Quick Reference]
\textbf{Core Question:} Which disciplines provide evidence for EBF, and what can each type of evidence be used for?

\textbf{The Three Disciplines:}
\begin{enumerate}[nosep]
    \item \textbf{Behavioral Economics (BE):} Calibrated parameters, causal identification
    \item \textbf{Behavioral Science (BS):} Mechanisms, theories, hypotheses
    \item \textbf{Neuroeconomics (NE):} Process evidence, dual-system support
\end{enumerate}

\textbf{Central Principle:} Different evidence types serve different EBF functions. Mixing them inappropriately leads to \textit{false precision}.

\textbf{Key Insight:} Only BE provides parameters for CORE-WHERE; BS and NE inform mechanisms but cannot calibrate the EBF equation directly.
\end{tcolorbox}

%===============================================================================
% ABSTRACT
%===============================================================================

\begin{abstract}
This appendix establishes the methodological boundaries between behavioral economics (BE), behavioral science (BS), and neuroeconomics (NE) for the Evidence-Based Framework (EBF). We argue that each discipline serves distinct functions: BE provides calibrated parameters and causal effects suitable for the EBF utility equation; BS provides mechanistic understanding and hypothesis generation; NE provides process evidence and dual-system validation. Conflating these roles leads to ``false precision''---treating qualitative insights as if they were quantitatively calibrated. We propose a classification system for the master bibliography and establish criteria for parameter validity. The appendix protects EBF's methodological integrity by ensuring that only appropriately validated evidence enters quantitative predictions.

\textbf{Important clarification:} This classification is \textit{not} a value hierarchy. Behavioral science provides essential contributions that behavioral economics cannot---including mechanistic understanding, hypothesis generation, and theoretical innovation. The classification serves only to clarify which evidence type serves which EBF function.
\end{abstract}

%===============================================================================
% POSITIONING STATEMENT
%===============================================================================

\begin{tcolorbox}[colback=purple!5!white, colframe=purple!75!black,
    title=\textbf{Positioning Statement: What This Appendix Is and Is Not}]

\textbf{This appendix is NOT:}
\begin{itemize}[nosep]
    \item A claim that behavioral economics is ``better'' than behavioral science
    \item A devaluation of psychological research
    \item An argument that qualitative insights are worthless
    \item Economics imperialism
\end{itemize}

\textbf{This appendix IS:}
\begin{itemize}[nosep]
    \item A \textbf{classification} of evidence types by function
    \item A \textbf{justification} of which scientific findings EBF builds upon for which purposes
    \item A \textbf{protection} against methodological errors (false precision)
    \item A \textbf{recognition} that different questions require different evidence
\end{itemize}

\textbf{The fundamental insight:}
\begin{quote}
\textit{EBF needs both behavioral economics AND behavioral science---but for different purposes. BE provides calibrated parameters for prediction ($\lambda = 2.25$). BS provides mechanistic understanding for intervention design (``why loss aversion exists''). Neither can replace the other. This appendix clarifies which evidence serves which function.}
\end{quote}

\end{tcolorbox}

%===============================================================================
% FUNDAMENTAL QUESTION
%===============================================================================

% -----------------------------------------------------------------------------
% APPENDIX SCOPE BOX
% -----------------------------------------------------------------------------

\begin{tcolorbox}[
    colback=orange!5!white,
    colframe=orange!75!black,
    title={\textbf{Appendix Scope: Ziel / In-Scope / Out-of-Scope / Constraints}},
    fonttitle=\bfseries
]

\textbf{Ziel (Objective):}
\begin{itemize}[nosep]
    \item The Problem:
\end{itemize}

\textbf{In-Scope:}
\begin{itemize}[nosep]
    \item Fundamental Question
    \item The Three Disciplines
    \item Evidence Classification for EBF
    \item The False Precision Problem
    \item The Problem of Implicit Evidence Hierarchies
\end{itemize}

\textbf{Out-of-Scope (delegiert an andere Appendices):}
\begin{itemize}[nosep]
    \item Glossary $\rightarrow$ Appendix G
    \item REF-GLOSSARY $\rightarrow$ Appendix G
    \item PREDICT-HABIT $\rightarrow$ Appendix HF
    \item LIT-LEARNING $\rightarrow$ Appendix XIX
    \item LIT-FEHR $\rightarrow$ Appendix K
\end{itemize}

\textbf{Constraints (Anwendungsgrenzen):}
\begin{itemize}[nosep]
    \item [TODO: Define when this appendix does NOT apply]
    \item [TODO: Define required prerequisites]
    \item [TODO: Define known limitations]
\end{itemize}

\textbf{Lieferobjekte (Deliverables):}
\begin{itemize}[nosep]
    \item 20 Table(s)
    \item Worked Example(s)
\end{itemize}

\end{tcolorbox}


\section{Fundamental Question}
\label{sec:evd-question}

\begin{tcolorbox}[colback=orange!5!white, colframe=orange!75!black,
    title=The Evidence Question]

\textbf{The Problem:}
\begin{quote}
EBF integrates insights from multiple disciplines. But not all evidence is created equal. A psychology paper showing ``first impressions matter'' is not the same as an economics paper estimating $\theta_1 = 0.35$ with IV identification.
\end{quote}

\textbf{The Risk:}
\begin{quote}
\textit{False Precision}---using qualitative or correlational findings as if they were causally identified parameters. This undermines prediction validity and intervention design.
\end{quote}

\textbf{The Question:}
\begin{quote}
\textit{What types of evidence can EBF use, and for what purposes?}
\end{quote}

\end{tcolorbox}


%%%%%%%%%%%%%%%%%%%%%%%%%%%%%%%%%%%%%%%%%%%%%%%%%%%%%%%%%%%%%%%%%%%%%%%%%%%%%%%
\section{The Three Disciplines}
\label{sec:evd-disciplines}
%%%%%%%%%%%%%%%%%%%%%%%%%%%%%%%%%%%%%%%%%%%%%%%%%%%%%%%%%%%%%%%%%%%%%%%%%%%%%%%

\subsection{Behavioral Economics (Verhaltensökonomie)}

\begin{tcolorbox}[colback=green!5!white, colframe=green!75!black,
    title=Definition: Behavioral Economics]

\textbf{Core Identity:} Economics with realistic psychological assumptions, maintaining the utility maximization framework.

\textbf{Key Characteristics:}
\begin{itemize}[nosep]
    \item Works \textit{within} expected utility / prospect theory framework
    \item Produces \textbf{calibrated parameters} ($\lambda$, $\beta$, $\theta$, $\gamma$)
    \item Uses \textbf{causal identification} (RCTs, IV, RDD, DiD)
    \item Incentive-compatible experimental designs
    \item Pre-registration increasingly standard
\end{itemize}

\textbf{Canonical Figures:}
Kahneman, Tversky, Thaler, Fehr, Camerer (experimental work), Laibson, O'Donoghue

\textbf{Canonical Journals:}
AER, QJE, JPE, Econometrica, REStud, JEEA

\textbf{EBF Role:} \textbf{Primary source for CORE-WHERE parameters}

\end{tcolorbox}

\subsection{Behavioral Science (Verhaltenswissenschaften)}

\begin{tcolorbox}[colback=blue!5!white, colframe=blue!75!black,
    title=Definition: Behavioral Science]

\textbf{Core Identity:} Interdisciplinary study of behavior, often without utility framework.

\textbf{Key Characteristics:}
\begin{itemize}[nosep]
    \item May reject utility maximization entirely
    \item Produces \textbf{theories and mechanisms} (not calibrated parameters)
    \item Often \textbf{correlational} or \textbf{qualitative}
    \item Constructs like ``Grit,'' ``Growth Mindset,'' ``Authenticity''
    \item Variable replication rates
\end{itemize}

\textbf{Canonical Figures:}
Dweck, Duckworth, Baumeister, Cialdini (persuasion), Haidt (moral psychology)

\textbf{Canonical Journals:}
Psychological Science, JPSP, Psychological Bulletin, PNAS (psychology)

\textbf{EBF Role:} \textbf{Hypothesis generation and mechanism understanding}

\end{tcolorbox}

\subsection{Neuroeconomics (Neuroökonomie)}

\begin{tcolorbox}[colback=purple!5!white, colframe=purple!75!black,
    title=Definition: Neuroeconomics]

\textbf{Core Identity:} Neural basis of economic decision-making.

\textbf{Key Characteristics:}
\begin{itemize}[nosep]
    \item Uses fMRI, EEG, lesion studies, pharmacological interventions
    \item Provides \textbf{process evidence} (what happens during decision?)
    \item Supports \textbf{dual-system theories} (System 1 vs 2)
    \item Often correlational (brain activation $\leftrightarrow$ choice)
    \item Rarely produces population-representative parameters
\end{itemize}

\textbf{Canonical Figures:}
Camerer (neuro work), Glimcher, Rangel, Fehr (social neuroscience), Loewenstein

\textbf{Canonical Journals:}
Neuron, Nature Neuroscience, Journal of Neuroscience, SCAN

\textbf{EBF Role:} \textbf{Process validation and mechanism support}

\end{tcolorbox}


%%%%%%%%%%%%%%%%%%%%%%%%%%%%%%%%%%%%%%%%%%%%%%%%%%%%%%%%%%%%%%%%%%%%%%%%%%%%%%%
\section{Evidence Classification for EBF}
\label{sec:evd-classification}
%%%%%%%%%%%%%%%%%%%%%%%%%%%%%%%%%%%%%%%%%%%%%%%%%%%%%%%%%%%%%%%%%%%%%%%%%%%%%%%

\subsection{The Evidence Hierarchy}

\begin{table}[htbp]
\centering
\caption{Evidence Types and EBF Uses}
\label{tab:evd-hierarchy}
\renewcommand{\arraystretch}{1.4}
\begin{tabular}{@{}p{2.5cm}p{3.5cm}p{3cm}p{3cm}@{}}
\toprule
\textbf{Evidence Type} & \textbf{Characteristics} & \textbf{EBF Use} & \textbf{Example} \\
\midrule
\textbf{Tier 1: Causal-Calibrated} & RCT/IV + parameter estimate & CORE-WHERE, PREDICT & $\theta_1 = 0.35$ (Jones et al.) \\
\textbf{Tier 2: Causal-Qualitative} & RCT without calibration & Mechanism validation & ``Incentives increase participation'' \\
\textbf{Tier 3: Correlational} & Observational association & Hypothesis generation & ``Optimists live longer'' \\
\textbf{Tier 4: Theoretical} & Model without data & Framework development & Dual-system theory \\
\textbf{Tier 5: Neural} & Brain data only & Process understanding & ``Insula activates for losses'' \\
\bottomrule
\end{tabular}
\end{table}

\subsection{What Each Tier Can Do}

\begin{tcolorbox}[colback=green!5!white, colframe=green!75!black,
    title=Tier 1: Causal-Calibrated (Gold Standard)]

\textbf{Can be used for:}
\begin{itemize}[nosep]
    \item Direct parameter entry in CORE-WHERE (BBB)
    \item Quantitative predictions in PREDICT appendices
    \item Intervention ROI calculations
    \item EBF equation calibration
\end{itemize}

\textbf{Requirements:}
\begin{itemize}[nosep]
    \item Causal identification (experimental or quasi-experimental)
    \item Numerical parameter estimate with standard errors
    \item Incentive-compatible design
    \item Replicable methodology
\end{itemize}

\textbf{Examples:}
\begin{itemize}[nosep]
    \item $\lambda = 2.25$ for loss aversion (Kahneman \& Tversky)
    \item $\theta_1 = 0.32\text{--}0.36$ for first-exposure effect (Jones et al. 2025)
    \item $\beta = 0.70$ for present bias (Laibson estimates)
\end{itemize}

\end{tcolorbox}

\begin{tcolorbox}[colback=blue!5!white, colframe=blue!75!black,
    title=Tier 2--4: Non-Calibrated Evidence]

\textbf{Can be used for:}
\begin{itemize}[nosep]
    \item Identifying \textit{which} mechanisms to model
    \item Generating hypotheses for PREDICT appendices
    \item Understanding \textit{why} parameters take certain values
    \item Qualitative validation of model structure
\end{itemize}

\textbf{Cannot be used for:}
\begin{itemize}[nosep]
    \item Direct parameter entry in CORE-WHERE
    \item Quantitative ROI predictions
    \item Precise intervention sizing
\end{itemize}

\textbf{Examples:}
\begin{itemize}[nosep]
    \item ``Growth mindset improves performance'' (Tier 3: correlational)
    \item ``People use mental accounting'' (Tier 4: theoretical)
    \item ``First impressions are important'' (Tier 3 without calibration)
\end{itemize}

\end{tcolorbox}

\begin{tcolorbox}[colback=purple!5!white, colframe=purple!75!black,
    title=Tier 5: Neural Evidence]

\textbf{Can be used for:}
\begin{itemize}[nosep]
    \item Validating that proposed mechanisms are neurally plausible
    \item Supporting dual-system (System 1/2) model structure
    \item Understanding \textit{process} of decision-making
    \item Ruling out implausible mechanisms
\end{itemize}

\textbf{Cannot be used for:}
\begin{itemize}[nosep]
    \item Direct parameter calibration
    \item Population-representative estimates
    \item Intervention design (fMRI $\neq$ behavior change)
\end{itemize}

\textbf{Special Case: Neuro-BE}
\begin{quote}
When neuroeconomics combines brain imaging with \textit{incentivized choices} and provides \textit{behavioral parameters}, it can qualify as Tier 1. Example: Camerer studies that estimate risk aversion from choices \textit{while} measuring brain activation.
\end{quote}

\end{tcolorbox}


%%%%%%%%%%%%%%%%%%%%%%%%%%%%%%%%%%%%%%%%%%%%%%%%%%%%%%%%%%%%%%%%%%%%%%%%%%%%%%%
\section{The False Precision Problem}
\label{sec:evd-false-precision}
%%%%%%%%%%%%%%%%%%%%%%%%%%%%%%%%%%%%%%%%%%%%%%%%%%%%%%%%%%%%%%%%%%%%%%%%%%%%%%%

\begin{tcolorbox}[colback=red!5!white, colframe=red!75!black,
    title=Warning: False Precision]

\textbf{Definition:} False precision occurs when qualitative or correlational findings are treated as if they were causally identified, calibrated parameters.

\textbf{Example of False Precision:}
\begin{quote}
``Dweck (2006) shows that growth mindset improves performance. Therefore, we estimate that a growth mindset intervention will increase productivity by 15\%.''
\end{quote}

\textbf{Why This Is Wrong:}
\begin{itemize}[nosep]
    \item Dweck's studies are largely correlational (Tier 3)
    \item No causal identification of mindset $\to$ performance
    \item ``15\%'' has no empirical basis
    \item Replication attempts show mixed results
\end{itemize}

\textbf{Correct Approach:}
\begin{quote}
``Dweck (2006) suggests that beliefs about ability may affect performance. This motivates \textit{hypothesis generation} for PREDICT, but parameter calibration requires Tier 1 evidence from incentivized experiments.''
\end{quote}

\end{tcolorbox}

\subsection{How False Precision Enters EBF}

\begin{enumerate}
\item \textbf{Mechanism $\to$ Parameter Leap:}
    \begin{quote}
    ``The literature shows X mechanism exists'' $\not\Rightarrow$ ``X has effect size Y''
    \end{quote}

\item \textbf{Lab $\to$ Field Extrapolation:}
    \begin{quote}
    Psychology lab effect $\not\Rightarrow$ Same effect in field with real stakes
    \end{quote}

\item \textbf{WEIRD $\to$ Universal Assumption:}
    \begin{quote}
    Effect in Western undergraduates $\not\Rightarrow$ Effect in DACH professionals
    \end{quote}

\item \textbf{Correlation $\to$ Causation:}
    \begin{quote}
    ``X correlates with Y'' $\not\Rightarrow$ ``Intervention on X causes Y''
    \end{quote}
\end{enumerate}


%%%%%%%%%%%%%%%%%%%%%%%%%%%%%%%%%%%%%%%%%%%%%%%%%%%%%%%%%%%%%%%%%%%%%%%%%%%%%%%
\section{The Problem of Implicit Evidence Hierarchies}
\label{sec:evd-implicit}
%%%%%%%%%%%%%%%%%%%%%%%%%%%%%%%%%%%%%%%%%%%%%%%%%%%%%%%%%%%%%%%%%%%%%%%%%%%%%%%

The false precision problem emerges from \textbf{implicit} rather than \textbf{explicit} evidence hierarchies. When researchers do not articulate their assumptions about evidence quality, systematic errors become invisible and uncorrectable \citep{ioannidis2005most}. This insight builds on decades of work in philosophy of science showing that much scientific knowledge is ``tacit''---we can know more than we can tell \citep{polanyi1966tacit}.

\subsection{Why Implicit Assumptions Are Dangerous}

\begin{tcolorbox}[colback=red!5!white, colframe=red!75!black,
    title=The Five Dangers of Implicit Evidence Hierarchies]

\textbf{1. Cherry-Picking Across Disciplines}
\begin{quote}
Without explicit tiers, researchers unconsciously select evidence that supports their conclusions---using BE parameters when convenient, BS mechanisms when useful, without acknowledging the category shift.
\end{quote}

\textbf{2. Unconscious Mixing of Evidence Types}
\begin{quote}
``Kahneman shows loss aversion ($\lambda = 2.25$), and Dweck shows growth mindset matters, so we model both...'' This sentence treats Tier 1 and Tier 3 evidence identically.
\end{quote}

\textbf{3. No Audit Trail for Parameter Choices}
\begin{quote}
When asked ``Why did you use $\theta_1 = 0.35$?'', the answer should be ``Jones et al. (2025), Tier 1, IV-identified''---not ``It seemed reasonable based on the literature.''
\end{quote}

\textbf{4. Inconsistent Standards Within Projects}
\begin{quote}
The same model uses Tier 1 evidence for $\lambda$ (loss aversion) but accepts Tier 3 correlational findings for ``motivation effects''---without noting the asymmetry.
\end{quote}

\textbf{5. Replication Illusion} \citep{simmons2011false}
\begin{quote}
``Research shows...'' conceals whether the source was an RCT with real stakes or a convenience sample survey. Both get cited as ``evidence.''
\end{quote}

\end{tcolorbox}

\subsection{The Hidden Cost of Implicit Hierarchies}

\begin{table}[htbp]
\centering
\caption{Implicit vs. Explicit Evidence Use: A Comparison}
\label{tab:evd-implicit-explicit}
\renewcommand{\arraystretch}{1.3}
\begin{tabular}{@{}p{3.5cm}p{5cm}p{5cm}@{}}
\toprule
\textbf{Aspect} & \textbf{Implicit Hierarchy} & \textbf{Explicit Hierarchy (EBF)} \\
\midrule
Parameter choice & Ad-hoc, based on familiarity & Tier-constrained, documented \\
Source quality & Invisible to reader/user & Tagged in BibTeX metadata \\
Audit trail & None---``the literature says'' & Full provenance: Paper $\to$ Tier $\to$ Parameter \\
Cross-project consistency & Variable, researcher-dependent & Standardized via METHOD-EVIDENCE \\
LLM behavior & Hallucination-prone, equal treatment & Constrained to Tier 1--2 for parameters \\
Error detection & Difficult---assumptions hidden & Easy---violations flagged automatically \\
Replication & Cannot verify evidence quality & Tier tags enable quality-matched replication \\
\bottomrule
\end{tabular}
\end{table}

\subsection{The Institutional Blindspot}

\begin{tcolorbox}[colback=orange!5!white, colframe=orange!75!black,
    title=Why Academia Tolerates Implicit Hierarchies]

\textbf{Within disciplines, implicit hierarchies work} \citep{kuhn1962structure, collins1985changing}:
\begin{itemize}[nosep]
    \item Economists know which journals require causal identification
    \item Psychologists know which findings replicate
    \item Shared tacit knowledge enables quality judgments \citep{polanyi1966tacit}
\end{itemize}

\textbf{Across disciplines, implicit hierarchies fail:}
\begin{itemize}[nosep]
    \item Economists may not know psychology's replication crisis \citep{collaboration2015estimating, camerer2016evaluating}
    \item Psychologists may not distinguish IV from OLS
    \item Practitioners (consultants, policymakers) lack disciplinary context
    \item LLMs have no tacit knowledge at all \citep{collins2010tacit}
\end{itemize}

\textbf{Key insight:}
\begin{quote}
\textit{Implicit evidence hierarchies rely on disciplinary socialization that does not transfer across fields---and cannot be encoded into computational systems. The EBF's explicit tier system makes tacit quality judgments visible and portable.}
\end{quote}

\end{tcolorbox}

\subsection{Case Study: What Goes Wrong Without Explicit Classification}

\begin{tcolorbox}[colback=yellow!5!white, colframe=yellow!75!black,
    title=Case Study: A Consulting Project Without Explicit Tiers]

\textbf{Scenario:} A consulting team designs a behavioral intervention to increase savings.

\textbf{What they cite (all treated equally):}
\begin{enumerate}[nosep]
    \item Thaler \& Benartzi (2004): Save More Tomorrow $\to$ 3.5pp increase
    \item Duckworth (2007): Grit predicts success $\to$ ``grit training''
    \item Ariely (2008): ``Predictably irrational'' $\to$ ``irrationality framing''
    \item Blog post: ``Gamification works'' $\to$ gamification elements
\end{enumerate}

\textbf{What they actually mixed:}
\begin{enumerate}[nosep]
    \item Tier 1: RCT with real retirement savings
    \item Tier 3: Correlational psychology, replication issues
    \item Tier 2: Popular science, some experiments
    \item Tier 4: Anecdote, no evidence
\end{enumerate}

\textbf{Result:} Intervention includes gamification and grit training (Tier 3--4) alongside defaults (Tier 1). When intervention underperforms, impossible to diagnose which component failed.

\textbf{With explicit tiers:} Team would have distinguished between ``proven mechanisms'' (Tier 1 defaults) and ``hypotheses to test'' (Tier 3--4 grit, gamification).

\end{tcolorbox}

\subsection{The Solution: Make the Hierarchy Explicit}

The solution follows the path pioneered by evidence-based medicine with the GRADE system \citep{guyatt2008grade} and the open science movement's push for transparency \citep{nosek2015promoting, munafo2017manifesto}:

\begin{enumerate}
\item \textbf{At the source level:} Tag every bibliography entry with \texttt{evidence\_tier}
\item \textbf{At the parameter level:} Document provenance in CORE-WHERE
\item \textbf{At the model level:} Distinguish ``calibrated'' from ``assumed'' parameters
\item \textbf{At the report level:} State evidence quality for each claim
\item \textbf{At the LLM level:} Constrain retrieval by tier
\end{enumerate}

\textbf{The fundamental principle:}
\begin{quote}
\textit{Every parameter in an EBF model must have explicit provenance: Source $\to$ Design $\to$ Tier $\to$ Confidence. Implicit assumptions are methodological debt that compounds over time.}
\end{quote}


%%%%%%%%%%%%%%%%%%%%%%%%%%%%%%%%%%%%%%%%%%%%%%%%%%%%%%%%%%%%%%%%%%%%%%%%%%%%%%%
\section{Bibliography Classification System}
\label{sec:evd-bibliography}
%%%%%%%%%%%%%%%%%%%%%%%%%%%%%%%%%%%%%%%%%%%%%%%%%%%%%%%%%%%%%%%%%%%%%%%%%%%%%%%

\subsection{Proposed Tags for bcm\_master.bib}

\begin{table}[htbp]
\centering
\caption{Bibliography Classification Tags}
\label{tab:evd-tags}
\renewcommand{\arraystretch}{1.3}
\begin{tabular}{@{}llp{6cm}@{}}
\toprule
\textbf{Tag} & \textbf{Meaning} & \textbf{Criteria} \\
\midrule
\texttt{be-experimental} & BE with causal ID & RCT, IV, RDD + incentives + parameter \\
\texttt{be-observational} & BE with observational data & Economic data, structural estimation \\
\texttt{bs-experimental} & Behavioral science RCT & Psychology RCT, often no incentives \\
\texttt{bs-correlational} & Behavioral science correlation & Survey, observational psychology \\
\texttt{neuro-behavioral} & Neuro + behavioral parameter & Brain imaging + incentivized choices \\
\texttt{neuro-mechanism} & Neuro for mechanism only & Brain imaging without behavior calibration \\
\texttt{theoretical} & Pure theory & Models, frameworks, no data \\
\bottomrule
\end{tabular}
\end{table}

\subsection{Implementation}

Add to each BibTeX entry:
\begin{verbatim}
@article{jones2025habits,
  ...
  keywords = {be-experimental, habit-formation, health},
  ebf-tier = {1},
  ...
}
\end{verbatim}


%%%%%%%%%%%%%%%%%%%%%%%%%%%%%%%%%%%%%%%%%%%%%%%%%%%%%%%%%%%%%%%%%%%%%%%%%%%%%%%
\section{Decision Rules for Evidence Use}
\label{sec:evd-rules}
%%%%%%%%%%%%%%%%%%%%%%%%%%%%%%%%%%%%%%%%%%%%%%%%%%%%%%%%%%%%%%%%%%%%%%%%%%%%%%%

\begin{tcolorbox}[colback=green!10!white, colframe=green!75!black,
    title=Evidence Use Decision Tree]

\textbf{Question 1: Does the study provide a numerical parameter estimate?}
\begin{itemize}[nosep]
    \item No $\Rightarrow$ Cannot use for CORE-WHERE; use for mechanism only
    \item Yes $\Rightarrow$ Continue to Question 2
\end{itemize}

\textbf{Question 2: Is the estimate causally identified?}
\begin{itemize}[nosep]
    \item No (correlational) $\Rightarrow$ Use with explicit uncertainty bounds; flag in BBB
    \item Yes $\Rightarrow$ Continue to Question 3
\end{itemize}

\textbf{Question 3: Was the design incentive-compatible?}
\begin{itemize}[nosep]
    \item No (hypothetical choices) $\Rightarrow$ Discount by factor $\delta \approx 0.5$
    \item Yes $\Rightarrow$ Continue to Question 4
\end{itemize}

\textbf{Question 4: Is the sample representative of target population?}
\begin{itemize}[nosep]
    \item No (WEIRD students) $\Rightarrow$ Note external validity limitation
    \item Yes $\Rightarrow$ \textbf{Tier 1: Use directly in CORE-WHERE}
\end{itemize}

\end{tcolorbox}


%%%%%%%%%%%%%%%%%%%%%%%%%%%%%%%%%%%%%%%%%%%%%%%%%%%%%%%%%%%%%%%%%%%%%%%%%%%%%%%
\section{Worked Example: Classifying Loss Aversion Evidence}
\label{sec:evd-worked-example}
%%%%%%%%%%%%%%%%%%%%%%%%%%%%%%%%%%%%%%%%%%%%%%%%%%%%%%%%%%%%%%%%%%%%%%%%%%%%%%%

\begin{tcolorbox}[colback=yellow!5!white, colframe=yellow!75!black,
    title=Worked Example: How to Classify Evidence for $\lambda$ (Loss Aversion)]

\textbf{Question:} We want to use loss aversion parameter $\lambda$ in an EBF model. What evidence exists, and how should we classify it?

\textbf{Step 1: Identify Candidate Sources}

\begin{table}[htbp]
\centering
\renewcommand{\arraystretch}{1.3}
\begin{tabular}{@{}p{4cm}p{3cm}p{2cm}p{2.5cm}@{}}
\toprule
\textbf{Source} & \textbf{Design} & \textbf{Parameter} & \textbf{Preliminary Tier} \\
\midrule
Kahneman \& Tversky (1992) & Experiment, incentives & $\lambda = 2.25$ & Tier 1? \\
Benartzi \& Thaler (1995) & Calibration to data & $\lambda = 2.5$ & Tier 2? \\
Tom et al. (2007) & fMRI + gambles & Neural correlate & Tier 5 \\
Baumeister et al. (2001) & Psychology review & ``Bad > Good'' & Tier 3 \\
\bottomrule
\end{tabular}
\end{table}

\textbf{Step 2: Apply Decision Tree}

\textit{Kahneman \& Tversky (1992):}
\begin{enumerate}[nosep]
\item Numerical parameter? \checkmark Yes ($\lambda = 2.25$)
\item Causally identified? \checkmark Yes (within-subject experimental)
\item Incentive-compatible? \checkmark Yes (real payoffs in some versions)
\item Representative sample? $\sim$ Mixed (students, but replicated widely)
\end{enumerate}
$\Rightarrow$ \textbf{Tier 1 with external validity note}

\textit{Tom et al. (2007):}
\begin{enumerate}[nosep]
\item Numerical parameter? $\times$ No (neural activation, not $\lambda$)
\end{enumerate}
$\Rightarrow$ \textbf{Tier 5: Process validation only}

\textit{Baumeister et al. (2001):}
\begin{enumerate}[nosep]
\item Numerical parameter? $\times$ No (qualitative: ``bad is stronger than good'')
\end{enumerate}
$\Rightarrow$ \textbf{Tier 3: Mechanism support only}

\textbf{Step 3: Decision for EBF}

\begin{itemize}[nosep]
\item \textbf{CORE-WHERE entry:} Use $\lambda = 2.25$ (K\&T 1992), cite as Tier 1
\item \textbf{Mechanism support:} Cite Baumeister for ``why'' loss aversion exists
\item \textbf{Neural validation:} Cite Tom et al. for process plausibility
\item \textbf{Do NOT:} Claim $\lambda = 2.25$ based on Baumeister (false precision)
\end{itemize}

\end{tcolorbox}

\subsection{Applying the Framework to First-Exposure Effect}

\begin{table}[htbp]
\centering
\caption{Evidence Classification for $\theta_1$ (First-Exposure Effect)}
\label{tab:evd-theta1-classification}
\renewcommand{\arraystretch}{1.3}
\begin{tabular}{@{}p{3.5cm}p{2.5cm}p{2cm}p{2cm}p{2cm}@{}}
\toprule
\textbf{Source} & \textbf{Design} & \textbf{Parameter} & \textbf{Tier} & \textbf{EBF Use} \\
\midrule
Jones et al. (2025) & RCT + IV & $\theta_1 = 0.35$ & \textbf{Tier 1} & CORE-WHERE \\
Ackerberg (2003) & Structural & Qualitative & Tier 2 & Mechanism \\
Dupas (2014) & RCT & Effect exists & Tier 2 & Validation \\
``First impressions matter'' & Folk wisdom & None & -- & Motivation \\
\bottomrule
\end{tabular}
\end{table}

\textbf{Correct statement:} ``Jones et al. (2025) provides Tier 1 evidence for $\theta_1 = 0.32$--$0.36$, supported by Tier 2 mechanism evidence from Ackerberg (2003) and Dupas (2014).''

\textbf{Incorrect statement:} ``Research shows first impressions matter, so we set $\theta_1 = 0.35$.'' (Skips evidence classification)


%%%%%%%%%%%%%%%%%%%%%%%%%%%%%%%%%%%%%%%%%%%%%%%%%%%%%%%%%%%%%%%%%%%%%%%%%%%%%%%
\section{Systematic Comparison: BE vs BS}
\label{sec:evd-comparison}
%%%%%%%%%%%%%%%%%%%%%%%%%%%%%%%%%%%%%%%%%%%%%%%%%%%%%%%%%%%%%%%%%%%%%%%%%%%%%%%

\subsection{The Eight Comparison Criteria}

\begin{table}[htbp]
\centering
\caption{Systematic Comparison of Behavioral Economics and Behavioral Science}
\label{tab:evd-be-bs-criteria}
\renewcommand{\arraystretch}{1.3}
\begin{tabular}{@{}p{3.5cm}p{5cm}p{5cm}@{}}
\toprule
\textbf{Criterion} & \textbf{Behavioral Economics} & \textbf{Behavioral Science} \\
\midrule
\multicolumn{3}{@{}l}{\textit{\textbf{Methodological Criteria}}} \\
1. Utility Framework & Works \textit{within} utility maximization & Often rejects or ignores utility \\
2. Parameter Output & Calibrated numbers ($\lambda$, $\beta$, $\theta$) & Qualitative mechanisms \\
3. Causal Identification & RCT, IV, DiD, RDD (explicit) & Often correlational \\
4. Incentive Compatibility & Real stakes (money, consequences) & Often hypothetical scenarios \\
5. Replication Rate & $\sim$60--70\% (Camerer et al. 2016) & $\sim$40\% (OSC 2015) \\
\addlinespace
\multicolumn{3}{@{}l}{\textit{\textbf{Institutional Criteria}}} \\
6. Publication Venues & AER, QJE, Econometrica, JEEA & Psych Science, JPSP, PNAS \\
7. Theoretical Connectivity & Welfare economics, mechanism design & Often isolated constructs \\
8. Falsifiability & Point predictions possible & Often only directional \\
\bottomrule
\end{tabular}
\end{table}

\subsection{The Three Critical Criteria for EBF}

For CORE-WHERE parameter eligibility, three criteria are decisive:

\begin{tcolorbox}[colback=red!5!white, colframe=red!75!black,
    title=Critical Criteria for Parameter Use]

\textbf{Criterion 1: Numerical Output?}
\begin{itemize}[nosep]
    \item BE: ``$\lambda = 2.25 \pm 0.3$'' $\Rightarrow$ \checkmark Usable for CORE-WHERE
    \item BS: ``Losses feel worse than gains'' $\Rightarrow$ $\times$ Not usable
\end{itemize}

\textbf{Criterion 2: Causally Identified?}
\begin{itemize}[nosep]
    \item BE: ``IV estimation shows $\theta_1 = 0.35$'' $\Rightarrow$ \checkmark Direct entry
    \item BS: ``Correlation $r = 0.4$'' $\Rightarrow$ $\times$ Hypothesis only
\end{itemize}

\textbf{Criterion 3: Incentive-Compatible?}
\begin{itemize}[nosep]
    \item BE: Real payoffs contingent on choices $\Rightarrow$ \checkmark Valid preferences
    \item BS: ``Imagine you won \$100...'' $\Rightarrow$ $\times$ Hypothetical bias
\end{itemize}

\end{tcolorbox}

\subsection{Publication Venue as Quality Signal}

\begin{table}[htbp]
\centering
\caption{Journal Tiers and Methodological Standards}
\label{tab:evd-journals}
\renewcommand{\arraystretch}{1.2}
\begin{tabular}{@{}lp{5cm}p{5cm}@{}}
\toprule
\textbf{Tier} & \textbf{Journals} & \textbf{Standards} \\
\midrule
\textbf{Tier A} & AER, QJE, JPE, Econometrica, REStud & Strict causal ID, replication expected, pre-registration increasing \\
\textbf{Tier B} & JEEA, AEJ journals, Management Science & Solid experimental, sometimes less identification \\
\textbf{Tier C} & Psychological Science, PNAS & Peer review, but documented replication problems \\
\textbf{Tier D} & JPSP, field-specific psychology & Variable standards, construct validity focus \\
\bottomrule
\end{tabular}
\end{table}

\textbf{Implication:} Publication venue is a \textit{proxy} for methodological rigor, not a guarantee. Individual paper assessment still required.

\subsection{Theoretical Connectivity (Anschlussfähigkeit)}

\begin{tcolorbox}[colback=blue!5!white, colframe=blue!75!black,
    title=Why Theoretical Connectivity Matters]

\textbf{BE findings are connectable to:}
\begin{itemize}[nosep]
    \item \textbf{Welfare Economics:} Can compute utility changes, deadweight loss
    \item \textbf{Mechanism Design:} Can optimize incentives given preferences
    \item \textbf{Policy Analysis:} Can conduct cost-benefit analysis
    \item \textbf{Prediction:} Can forecast behavior under counterfactuals
\end{itemize}

\textbf{BS findings are often NOT connectable:}
\begin{itemize}[nosep]
    \item ``Grit'' $\to$ No utility equivalent, no welfare function
    \item ``Growth Mindset'' $\to$ Cannot optimize, no counterfactual
    \item ``Authenticity'' $\to$ Not translatable to $u(x)$
\end{itemize}

\textbf{Key insight:} BE maintains \textit{theoretical continuity} with economics. BS often creates \textit{isolated constructs} that cannot be integrated into optimization frameworks.

\end{tcolorbox}


%%%%%%%%%%%%%%%%%%%%%%%%%%%%%%%%%%%%%%%%%%%%%%%%%%%%%%%%%%%%%%%%%%%%%%%%%%%%%%%
\section{The European Laboratory Standard}
\label{sec:evd-european-lab}
%%%%%%%%%%%%%%%%%%%%%%%%%%%%%%%%%%%%%%%%%%%%%%%%%%%%%%%%%%%%%%%%%%%%%%%%%%%%%%%

\subsection{Laboratory Economics: BE vs BS Methods}

Both BE and BS conduct laboratory experiments. The difference lies in \textit{how} they use the laboratory:

\begin{table}[htbp]
\centering
\caption{Laboratory Methods: BE vs BS}
\label{tab:evd-lab-comparison}
\renewcommand{\arraystretch}{1.3}
\begin{tabular}{@{}p{3.5cm}p{5cm}p{5cm}@{}}
\toprule
\textbf{Aspect} & \textbf{BE Laboratory} & \textbf{BS Laboratory} \\
\midrule
Payoffs & Real money (€5--50 per decision) & Hypothetical or course credit \\
Goal & Calibrate parameters ($\lambda = 2.25$) & Demonstrate effect ($p < 0.05$) \\
Design & Within-subject, randomization & Between-groups, convenience \\
Output & Structural estimation & Mean comparison \\
Protocol & Published, pre-registered & Variable \\
Replication & Explicit replication studies & Often absent \\
\bottomrule
\end{tabular}
\end{table}

\subsection{The European Laboratory Standard}

The European experimental economics tradition, centered at institutions like Zurich, Bonn, Innsbruck, Vienna, Amsterdam, and Nottingham, established rigorous methodological standards that distinguish BE laboratory research:

\begin{tcolorbox}[colback=green!5!white, colframe=green!75!black,
    title=European Laboratory Standard: Core Principles]

\textbf{Principle 1: Incentive Compatibility}
\begin{itemize}[nosep]
    \item All choices have real monetary consequences
    \item No deception (unlike some psychology traditions)
    \item Payoffs calibrated to ensure salience (typically €10--50)
\end{itemize}

\textbf{Principle 2: Anonymity and Control}
\begin{itemize}[nosep]
    \item Double-blind protocols where possible
    \item No reputation effects (one-shot or stranger matching)
    \item Computerized interaction to minimize experimenter effects
\end{itemize}

\textbf{Principle 3: Structural Estimation}
\begin{itemize}[nosep]
    \item Not just ``does X affect Y?'' but ``by how much?''
    \item Parameter recovery: estimate $\alpha$, $\beta$, $\lambda$ from choice data
    \item Model comparison via likelihood ratios
\end{itemize}

\textbf{Principle 4: Replicability}
\begin{itemize}[nosep]
    \item Full protocol publication (instructions, software)
    \item z-Tree or oTree code shared
    \item Explicit replication attempts valued
\end{itemize}

\end{tcolorbox}

\subsection{Canonical Paradigms}

\begin{table}[htbp]
\centering
\caption{European Laboratory Standard: Canonical Games}
\label{tab:evd-canonical-games}
\renewcommand{\arraystretch}{1.3}
\begin{tabular}{@{}p{3cm}p{3cm}p{4cm}p{3cm}@{}}
\toprule
\textbf{Game} & \textbf{Design} & \textbf{Parameter Output} & \textbf{Typical Stakes} \\
\midrule
Ultimatum Game & 2 players, Proposer/Responder & Rejection rate, fairness threshold & €10--40 \\
Dictator Game & 1 player allocates to passive other & Pure altruism parameter & €10--20 \\
Public Goods Game & N players, contribution decision & Cooperation rate, $\gamma$ (MPCR) & €10--30 \\
Trust Game & Sender/Receiver, sequential & Trust and reciprocity parameters & €15--50 \\
Gift Exchange & Employer/Worker, wage-effort & Reciprocity elasticity & €15--40 \\
Risk Elicitation & Lottery choices (Holt-Laury) & Risk aversion coefficient $r$ & €5--30 \\
Time Preference & Sooner-smaller vs later-larger & $\beta$, $\delta$ parameters & €10--100 \\
\bottomrule
\end{tabular}
\end{table}

\subsection{Example: Loss Aversion in BE vs BS Laboratory}

\begin{tcolorbox}[colback=yellow!5!white, colframe=yellow!75!black,
    title=Same Phenomenon, Different Laboratory Approaches]

\textbf{BE Laboratory (European Standard):}
\begin{quote}
``Choose between:\\
A) 50\% chance of €100, 50\% chance of losing €50\\
B) €0 for certain''\\
$\to$ Real payoff at end of session\\
$\to$ From choice patterns: estimate $\lambda = 2.25$\\
$\to$ Incentive-compatible, within-subject, structural
\end{quote}

\textbf{BS Laboratory (Typical):}
\begin{quote}
``Imagine you lost €100. How does this make you feel? (1--7 scale)''\\
``Now imagine you gained €100. How does this make you feel?''\\
$\to$ Compare mean ratings\\
$\to$ Conclude: ``Losses feel worse than gains'' ($p < 0.05$)\\
$\to$ No $\lambda$, no real stakes, no structural estimation
\end{quote}

\textbf{Key difference:} Both are laboratory experiments. But only BE extracts \textit{calibrated parameters} from \textit{incentive-compatible decisions}.

\end{tcolorbox}

\subsection{Why the European Standard Matters for EBF}

\begin{enumerate}
\item \textbf{Parameter Validity:} Only incentive-compatible elicitation reveals true preferences
\item \textbf{Quantification:} Structural estimation provides numbers for CORE-WHERE
\item \textbf{Replicability:} Published protocols enable verification
\item \textbf{External Validity:} Real stakes approximate field decisions better than hypotheticals
\end{enumerate}

\textbf{Implication for EBF:} Parameters from European Laboratory Standard research (Tier 1) can enter CORE-WHERE directly. Parameters from non-incentivized psychology laboratories require additional validation.


%%%%%%%%%%%%%%%%%%%%%%%%%%%%%%%%%%%%%%%%%%%%%%%%%%%%%%%%%%%%%%%%%%%%%%%%%%%%%%%
\section{The Fehr Methodological Foundation}
\label{sec:evd-fehr-foundation}
%%%%%%%%%%%%%%%%%%%%%%%%%%%%%%%%%%%%%%%%%%%%%%%%%%%%%%%%%%%%%%%%%%%%%%%%%%%%%%%

Ernst Fehr's work at the University of Zurich exemplifies and grounds the BE/BS distinction. His methodological approach provides the template for transforming psychological insights into calibrated economic parameters.

\subsection{The Fehr Transformation: BS $\to$ BE}

\begin{tcolorbox}[colback=green!5!white, colframe=green!75!black,
    title=The Fehr Recipe: From Concept to Parameter]

\textbf{Step 1: Identify a Psychological Phenomenon}
\begin{quote}
BS insight: ``People care about fairness and punish unfair behavior''
\end{quote}

\textbf{Step 2: Formalize in Utility Framework}
\begin{quote}
Fehr \& Schmidt (1999):
$$U_i(\pi) = \pi_i - \alpha_i \max\{\pi_j - \pi_i, 0\} - \beta_i \max\{\pi_i - \pi_j, 0\}$$
where $\alpha$ = disadvantageous inequity aversion, $\beta$ = advantageous inequity aversion
\end{quote}

\textbf{Step 3: Design Incentive-Compatible Game}
\begin{quote}
Ultimatum Game: Real money (CHF 20--40), one-shot, anonymous, computerized
\end{quote}

\textbf{Step 4: Estimate Parameters Structurally}
\begin{quote}
From rejection behavior: $\alpha \approx 0.5$--$1.0$, $\beta \approx 0.25$--$0.6$
\end{quote}

\textbf{Step 5: Replicate and Validate}
\begin{quote}
Cross-cultural replications, meta-analyses, field validation
\end{quote}

\end{tcolorbox}

\subsection{Foundational Papers: The BE/BS Transformation in Practice}

\begin{table}[htbp]
\centering
\caption{Fehr's Transformation of Psychological Concepts to Economic Parameters}
\label{tab:evd-fehr-papers}
\renewcommand{\arraystretch}{1.4}
\begin{tabular}{@{}p{3.5cm}p{3cm}p{4cm}p{3cm}@{}}
\toprule
\textbf{Paper} & \textbf{BS Concept} & \textbf{BE Formalization} & \textbf{Parameters} \\
\midrule
Fehr \& Schmidt (1999) \textit{QJE} & ``Fairness matters'' & Inequity aversion utility & $\alpha \approx 0.5$, $\beta \approx 0.25$ \\
\addlinespace
Fehr \& Gächter (2000) \textit{AER} & ``Free-riders get punished'' & Altruistic punishment function & Cost ratio $\sim$3:1 \\
\addlinespace
Fehr, Kirchsteiger \& Riedl (1993) \textit{QJE} & ``Workers reciprocate kindness'' & Gift exchange wage-effort & Reciprocity elasticity \\
\addlinespace
Fehr \& Fischbacher (2004) \textit{EE} & ``Third parties enforce norms'' & Third-party punishment & Punishment demand curve \\
\addlinespace
Fehr \& Camerer (2007) \textit{Science} & ``Social preferences exist'' & Neuroeconomic validation & Neural correlates of $\alpha$, $\beta$ \\
\bottomrule
\end{tabular}
\end{table}

\subsection{Why Fehr's Approach Grounds the Distinction}

\begin{tcolorbox}[colback=blue!5!white, colframe=blue!75!black,
    title=The Fehr Contribution to BE/BS Methodology]

\textbf{1. Demonstrated that BS concepts CAN be parameterized:}
\begin{itemize}[nosep]
    \item ``Fairness'' $\to$ $\alpha$, $\beta$ in Fehr-Schmidt model
    \item ``Trust'' $\to$ Sending rate in Trust Game
    \item ``Reciprocity'' $\to$ Return rate, punishment rate
    \item ``Cooperation'' $\to$ Contribution rate in Public Goods
\end{itemize}

\textbf{2. Established methodological standards that DISTINGUISH BE from BS:}
\begin{itemize}[nosep]
    \item No deception (unlike Milgram, Asch traditions)
    \item Real monetary stakes (unlike hypothetical scenarios)
    \item Structural estimation (unlike mean comparisons)
    \item One-shot designs (unlike repeated interaction confounds)
\end{itemize}

\textbf{3. Created replicable paradigms:}
\begin{itemize}[nosep]
    \item Ultimatum Game $\to$ standard fairness measure
    \item Trust Game $\to$ standard trust/reciprocity measure
    \item Public Goods Game $\to$ standard cooperation measure
    \item Gift Exchange $\to$ standard labor reciprocity measure
\end{itemize}

\end{tcolorbox}

\subsection{The Fehr Standard vs. Psychology Standard}

\begin{table}[htbp]
\centering
\caption{Methodological Comparison: Fehr Standard vs. Typical Psychology}
\label{tab:evd-fehr-vs-psych}
\renewcommand{\arraystretch}{1.3}
\begin{tabular}{@{}p{4cm}p{5cm}p{5cm}@{}}
\toprule
\textbf{Dimension} & \textbf{Fehr Standard (BE)} & \textbf{Typical Psychology (BS)} \\
\midrule
\textit{On ``Fairness''} & & \\
Question asked & ``How much will you sacrifice to punish unfairness?'' & ``How important is fairness to you? (1--7)'' \\
Stakes & CHF 20--40 real money & Hypothetical or course credit \\
Output & $\alpha = 0.5$, rejection threshold & ``Fairness is important'' ($p < 0.05$) \\
\addlinespace
\textit{On ``Trust''} & & \\
Question asked & ``How much of CHF 10 will you send to stranger?'' & ``Do you generally trust people? (1--7)'' \\
Stakes & Real transfer, tripled, returned or not & Self-report \\
Output & Trust rate = 0.52, reciprocity rate = 0.37 & Trust scale score \\
\addlinespace
\textit{On ``Cooperation''} & & \\
Question asked & ``How much will you contribute to public good?'' & ``Are you a cooperative person?'' \\
Stakes & Real payoff = $f$(own + others' contributions) & Self-assessment \\
Output & Contribution rate, decay function, punishment effect & Cooperation score \\
\bottomrule
\end{tabular}
\end{table}

\subsection{Implications for EBF Evidence Classification}

\begin{tcolorbox}[colback=yellow!5!white, colframe=yellow!75!black,
    title=The Fehr Test for Tier 1 Evidence]

When classifying evidence, ask: \textit{``Does this study meet the Fehr Standard?''}

\textbf{Fehr Standard Checklist:}
\begin{enumerate}[nosep]
    \item[$\square$] Real monetary stakes (not hypothetical)
    \item[$\square$] Incentive-compatible mechanism (choices have consequences)
    \item[$\square$] No deception (participants know payoff structure)
    \item[$\square$] Structural parameter estimation (not just significance test)
    \item[$\square$] Replicable design (protocol published)
    \item[$\square$] Anonymity preserved (no reputation effects)
\end{enumerate}

\textbf{If all boxes checked:} Tier 1 evidence, eligible for CORE-WHERE

\textbf{If some boxes unchecked:} Tier 2--3, mechanism support only

\end{tcolorbox}

\subsection{Fehr's Legacy for the BE/BS Distinction}

\begin{enumerate}
\item \textbf{Proof of concept:} Showed that ``soft'' psychological concepts (fairness, trust, reciprocity) can be made ``hard'' through proper methodology.

\item \textbf{Methodological template:} The Fehr Recipe (concept $\to$ utility $\to$ game $\to$ parameter $\to$ replicate) is the standard for BE research.

\item \textbf{Quality benchmark:} ``Does it meet the Fehr Standard?'' is a useful heuristic for evidence classification.

\item \textbf{Integration model:} Fehr's work shows how BS insights (fairness matters) become BE parameters ($\alpha$, $\beta$) without losing the underlying psychology.
\end{enumerate}

\textbf{Key insight:} The BE/BS distinction is not about \textit{rejecting} psychology---it's about \textit{transforming} psychological insights into quantifiable, falsifiable, policy-relevant parameters. Fehr's work is the paradigm case of this transformation.


%%%%%%%%%%%%%%%%%%%%%%%%%%%%%%%%%%%%%%%%%%%%%%%%%%%%%%%%%%%%%%%%%%%%%%%%%%%%%%%
\section{Historical Context: The Emergence of Behavioral Economics}
\label{sec:evd-history}
%%%%%%%%%%%%%%%%%%%%%%%%%%%%%%%%%%%%%%%%%%%%%%%%%%%%%%%%%%%%%%%%%%%%%%%%%%%%%%%

\subsection{Why BE and BS Diverged}

\begin{tcolorbox}[colback=gray!5!white, colframe=gray!75!black,
    title=Historical Timeline: From Anomalies to Discipline]

\textbf{1979: Prospect Theory (Kahneman \& Tversky)}
\begin{itemize}[nosep]
\item Psychology insight: People are loss averse
\item \textit{Critical move:} Formalized as utility function with $\lambda$ parameter
\item $\Rightarrow$ Made psychological insight \textit{economics-compatible}
\end{itemize}

\textbf{1980s: Anomalies Literature (Thaler)}
\begin{itemize}[nosep]
\item Documented systematic violations of standard theory
\item Mental accounting, endowment effect, fairness concerns
\item Published in \textit{Journal of Economic Perspectives} (economics outlet)
\end{itemize}

\textbf{1988: Rational Addiction (Becker \& Murphy)}
\begin{itemize}[nosep]
\item Chicago response: Keep utility framework, add dynamics
\item Showed ``irrational'' behavior can emerge from rational optimization
\item Set template for BE: Modify preferences, keep optimization
\end{itemize}

\textbf{1999--2003: Experimental Economics Boom}
\begin{itemize}[nosep]
\item Fehr \& Schmidt (1999): Inequity aversion with $\alpha, \beta$ parameters
\item Charness \& Rabin (2002): Social preferences taxonomy
\item \textit{Key:} All models retained utility maximization
\end{itemize}

\textbf{2002: Nobel Prize to Kahneman}
\begin{itemize}[nosep]
\item Legitimized behavioral approach within economics
\item Tversky (deceased) would have shared
\end{itemize}

\textbf{2017: Nobel Prize to Thaler}
\begin{itemize}[nosep]
\item BE fully established as economics subfield
\item Distinct from psychology: Same phenomena, different formalization
\end{itemize}

\end{tcolorbox}

\subsection{The Crucial Difference: Parameterization}

\begin{table}[htbp]
\centering
\caption{How BE and BS Treat the Same Phenomenon}
\label{tab:evd-be-bs-comparison}
\renewcommand{\arraystretch}{1.4}
\begin{tabular}{@{}p{3cm}p{5cm}p{5cm}@{}}
\toprule
\textbf{Phenomenon} & \textbf{Behavioral Science Treatment} & \textbf{Behavioral Economics Treatment} \\
\midrule
Loss aversion & ``Losses loom larger than gains'' (qualitative) & $U(x) = x$ if $x \geq 0$; $U(x) = \lambda |x|$ if $x < 0$, with $\lambda = 2.25$ \\
Present bias & ``People prefer immediate gratification'' & $U_0 = u_0 + \beta \sum_{t=1}^T \delta^t u_t$, with $\beta = 0.70$ \\
Social preferences & ``Fairness matters'' & $U_i = \pi_i - \alpha \max(\pi_j - \pi_i, 0) - \beta \max(\pi_i - \pi_j, 0)$ \\
First impressions & ``First impressions matter'' & $\theta_1 = 0.35$, $\theta_2 \approx 0$ (IV-identified) \\
\bottomrule
\end{tabular}
\end{table}

\textbf{Key insight:} BE doesn't just say ``X exists''---it says ``X has magnitude Y, identified via method Z.''


%%%%%%%%%%%%%%%%%%%%%%%%%%%%%%%%%%%%%%%%%%%%%%%%%%%%%%%%%%%%%%%%%%%%%%%%%%%%%%%
\section{LIT-Appendix Evidence Classification}
\label{sec:evd-lit-classification}
%%%%%%%%%%%%%%%%%%%%%%%%%%%%%%%%%%%%%%%%%%%%%%%%%%%%%%%%%%%%%%%%%%%%%%%%%%%%%%%

The EBF contains 33 LIT appendices integrating research by author or topic. Each has a dominant evidence tier:

\begin{table}[htbp]
\centering
\caption{Evidence Tier Classification of LIT Appendices}
\label{tab:evd-lit-tiers}
\renewcommand{\arraystretch}{1.2}
\begin{tabular}{@{}llcl@{}}
\toprule
\textbf{Code} & \textbf{LIT Appendix} & \textbf{Dominant Tier} & \textbf{Rationale} \\
\midrule
\multicolumn{4}{@{}l}{\textit{\textbf{Tier 1 Dominant (Experimental BE)}}} \\
K & LIT-FEHR & 1 & Lab experiments with incentives, $\alpha, \beta, \gamma$ \\
L & LIT-THALER & 1--2 & Field experiments, some calibration \\
N & LIT-KAHNEMAN & 1 & Prospect theory parameters \\
T & LIT-CAMERER & 1 & Experimental economics, neuro-behavioral \\
M & LIT-GNEEZY & 1 & Incentive experiments \\
XIX & LIT-LEARNING & 1 & IV-identified habit parameters \\
\addlinespace
\multicolumn{4}{@{}l}{\textit{\textbf{Tier 2 Dominant (Causal, Qualitative)}}} \\
O & LIT-ARIELY & 2 & Experiments, often qualitative effects \\
I & LIT-DECI & 2--3 & Motivation research, mixed methods \\
P & LIT-AKERLOF & 2 & Theory + structural estimation \\
Q & LIT-LOEWENSTEIN & 2 & Visceral factors, some experiments \\
\addlinespace
\multicolumn{4}{@{}l}{\textit{\textbf{Tier 3--4 Dominant (Correlational/Theoretical)}}} \\
U & LIT-METAANALYSIS & 3 & Aggregates mixed-quality studies \\
J & LIT-DUCKWORTH & 3 & Grit research, correlational \\
-- & LIT-DWECK & 3 & Growth mindset, replication issues \\
\addlinespace
\multicolumn{4}{@{}l}{\textit{\textbf{Tier 5 (Neural)}}} \\
-- & (Neuro sections in T) & 5 & Process evidence only \\
\bottomrule
\end{tabular}
\end{table}

\subsection{Implications for EBF Parameter Use}

\begin{tcolorbox}[colback=green!5!white, colframe=green!75!black,
    title=Using LIT Appendices for CORE-WHERE]

\textbf{Direct parameter extraction (Tier 1 LIT):}
\begin{itemize}[nosep]
\item K (LIT-FEHR): $\alpha$, $\beta$ for inequity aversion
\item N (LIT-KAHNEMAN): $\lambda$ for loss aversion, $\pi(p)$ for probability weighting
\item XIX (LIT-LEARNING): $\theta_1$, $\theta_2$ for habit formation
\end{itemize}

\textbf{Mechanism understanding (Tier 2--3 LIT):}
\begin{itemize}[nosep]
\item I (LIT-DECI): Why intrinsic motivation matters (but not how much)
\item Q (LIT-LOEWENSTEIN): When visceral factors dominate (but not calibrated)
\end{itemize}

\textbf{Hypothesis generation (Tier 3--4 LIT):}
\begin{itemize}[nosep]
\item J (LIT-DUCKWORTH): Persistence may matter---test with Tier 1 design
\item Dweck: Beliefs about ability may matter---requires BE validation
\end{itemize}

\end{tcolorbox}


%%%%%%%%%%%%%%%%%%%%%%%%%%%%%%%%%%%%%%%%%%%%%%%%%%%%%%%%%%%%%%%%%%%%%%%%%%%%%%%
\section{CORE-WHERE Parameter Provenance}
\label{sec:evd-bbb-mapping}
%%%%%%%%%%%%%%%%%%%%%%%%%%%%%%%%%%%%%%%%%%%%%%%%%%%%%%%%%%%%%%%%%%%%%%%%%%%%%%%

\begin{table}[htbp]
\centering
\caption{Evidence Tiers for Key BBB Parameters}
\label{tab:evd-bbb-tiers}
\renewcommand{\arraystretch}{1.3}
\begin{tabular}{@{}p{2.5cm}cp{3.5cm}cp{3cm}@{}}
\toprule
\textbf{Parameter} & \textbf{Value} & \textbf{Source} & \textbf{Tier} & \textbf{Confidence} \\
\midrule
\multicolumn{5}{@{}l}{\textit{\textbf{Prospect Theory Parameters}}} \\
$\lambda$ (loss aversion) & 2.25 & Kahneman \& Tversky (1992) & 1 & High \\
$\alpha$ (diminishing sensitivity) & 0.88 & Kahneman \& Tversky (1992) & 1 & High \\
$\gamma$ (prob. weighting) & 0.61 & Prelec (1998) & 1 & Medium \\
\addlinespace
\multicolumn{5}{@{}l}{\textit{\textbf{Time Preference Parameters}}} \\
$\beta$ (present bias) & 0.70 & Laibson (1997) & 1 & High \\
$\delta$ (discount factor) & 0.95--0.99 & Multiple & 1 & High \\
\addlinespace
\multicolumn{5}{@{}l}{\textit{\textbf{Social Preference Parameters}}} \\
$\alpha$ (disadvantageous ineq.) & 0.5--1.0 & Fehr \& Schmidt (1999) & 1 & Medium \\
$\beta$ (advantageous ineq.) & 0.25--0.6 & Fehr \& Schmidt (1999) & 1 & Medium \\
\addlinespace
\multicolumn{5}{@{}l}{\textit{\textbf{Habit Formation Parameters}}} \\
$\theta_1$ (first exposure) & 0.32--0.36 & Jones et al. (2025) & 1 & High \\
$\theta_2$ (second exposure) & $\approx 0.03$ & Jones et al. (2025) & 1 & High \\
\addlinespace
\multicolumn{5}{@{}l}{\textit{\textbf{Parameters Requiring Caution}}} \\
``Grit'' effect & -- & Duckworth et al. & 3 & Low \\
``Growth mindset'' & -- & Dweck et al. & 3 & Low \\
``Ego depletion'' & -- & Baumeister et al. & 3 & Very Low \\
\bottomrule
\end{tabular}
\end{table}

\textbf{Note:} Parameters marked Tier 3 with ``Low'' confidence should not be used for quantitative CORE-WHERE entries. They may inform qualitative mechanism discussions only.


%%%%%%%%%%%%%%%%%%%%%%%%%%%%%%%%%%%%%%%%%%%%%%%%%%%%%%%%%%%%%%%%%%%%%%%%%%%%%%%
\section{Neuroeconomics: A Detailed Treatment}
\label{sec:evd-neuro}
%%%%%%%%%%%%%%%%%%%%%%%%%%%%%%%%%%%%%%%%%%%%%%%%%%%%%%%%%%%%%%%%%%%%%%%%%%%%%%%

\subsection{What Neuroeconomics Can Contribute}

\begin{table}[htbp]
\centering
\caption{Neuroeconomics Contributions to EBF}
\label{tab:evd-neuro-contributions}
\renewcommand{\arraystretch}{1.3}
\begin{tabular}{@{}p{4cm}p{4cm}p{4cm}@{}}
\toprule
\textbf{Contribution} & \textbf{Example} & \textbf{EBF Use} \\
\midrule
Dual-system validation & vmPFC for value, dlPFC for control & Supports CORE-AWARE structure \\
Process timing & Decisions before conscious awareness & Informs System 1/2 boundary \\
Mechanism localization & Insula for loss, striatum for gain & Validates loss aversion as real \\
Individual differences & Neural predictors of risk tolerance & Informs segmentation (CORE-WHO) \\
\bottomrule
\end{tabular}
\end{table}

\subsection{What Neuroeconomics Cannot Provide}

\begin{tcolorbox}[colback=red!5!white, colframe=red!75!black,
    title=Limitations of Neural Evidence for EBF]

\textbf{1. Population Parameters:}
\begin{quote}
fMRI studies have N $\approx$ 20--50. Cannot estimate population-representative $\lambda$ or $\beta$.
\end{quote}

\textbf{2. Causal Intervention Effects:}
\begin{quote}
``Insula activates for losses'' $\neq$ ``Reducing insula activation reduces loss aversion''
\end{quote}

\textbf{3. Field Validity:}
\begin{quote}
Scanner environment $\neq$ Real-world decision context
\end{quote}

\textbf{4. Utility Parameters:}
\begin{quote}
Brain activation is not utility. BOLD signal $\not\Rightarrow$ $u(x)$
\end{quote}

\end{tcolorbox}

\subsection{The Neuro-BE Bridge}

Some neuroeconomics studies \textit{do} qualify for Tier 1 if they:
\begin{enumerate}[nosep]
    \item Use incentive-compatible behavioral tasks
    \item Estimate behavioral parameters (not just brain parameters)
    \item Have adequate sample sizes
    \item Report both neural and behavioral results
\end{enumerate}

\textbf{Example:} Camerer \& Fehr (2013) combines neural measurement with behavioral parameter estimation from ultimatum games.


%%%%%%%%%%%%%%%%%%%%%%%%%%%%%%%%%%%%%%%%%%%%%%%%%%%%%%%%%%%%%%%%%%%%%%%%%%%%%%%
\section{Results: Discipline Mapping}
\label{sec:evd-results}
%%%%%%%%%%%%%%%%%%%%%%%%%%%%%%%%%%%%%%%%%%%%%%%%%%%%%%%%%%%%%%%%%%%%%%%%%%%%%%%

\begin{tcolorbox}[colback=purple!5!white, colframe=purple!75!black,
    title=Summary: What Each Discipline Provides]

\begin{table}[htbp]
\centering
\renewcommand{\arraystretch}{1.4}
\begin{tabular}{@{}lccc@{}}
\toprule
\textbf{EBF Component} & \textbf{BE} & \textbf{BS} & \textbf{NE} \\
\midrule
CORE-WHERE parameters & \checkmark\checkmark\checkmark & $\times$ & $\circ$ \\
PREDICT calibration & \checkmark\checkmark\checkmark & $\times$ & $\times$ \\
Mechanism identification & \checkmark\checkmark & \checkmark\checkmark & \checkmark\checkmark \\
Hypothesis generation & \checkmark & \checkmark\checkmark\checkmark & \checkmark \\
Process understanding & \checkmark & \checkmark & \checkmark\checkmark\checkmark \\
Dual-system support & \checkmark & \checkmark & \checkmark\checkmark\checkmark \\
Intervention design & \checkmark\checkmark\checkmark & \checkmark & $\times$ \\
\bottomrule
\end{tabular}
\end{table}

\textbf{Legend:} \checkmark\checkmark\checkmark = Primary source; \checkmark\checkmark = Useful; \checkmark = Limited; $\circ$ = Conditional; $\times$ = Not appropriate

\end{tcolorbox}


%%%%%%%%%%%%%%%%%%%%%%%%%%%%%%%%%%%%%%%%%%%%%%%%%%%%%%%%%%%%%%%%%%%%%%%%%%%%%%%
\section{Glossary of Terms}
\label{sec:evd-glossary}
%%%%%%%%%%%%%%%%%%%%%%%%%%%%%%%%%%%%%%%%%%%%%%%%%%%%%%%%%%%%%%%%%%%%%%%%%%%%%%%

\textit{For complete definitions, see Appendix G (REF-GLOSSARY).}

\begin{table}[htbp]
\centering
\caption{Methodological Terms}
\label{tab:evd-glossary}
\renewcommand{\arraystretch}{1.3}
\begin{tabular}{@{}lp{8cm}c@{}}
\toprule
\textbf{Term} & \textbf{Definition} & \textbf{In G?} \\
\midrule
Causal identification & Design that permits causal inference (not just correlation) & NEW \\
Incentive compatibility & Participants have real stakes in their choices & \checkmark \\
False precision & Treating qualitative findings as calibrated parameters & NEW \\
WEIRD & Western, Educated, Industrialized, Rich, Democratic samples & NEW \\
External validity & Generalizability beyond study sample & \checkmark \\
Parameter calibration & Estimating numerical values for model parameters & NEW \\
\bottomrule
\end{tabular}
\end{table}


%%%%%%%%%%%%%%%%%%%%%%%%%%%%%%%%%%%%%%%%%%%%%%%%%%%%%%%%%%%%%%%%%%%%%%%%%%%%%%%
\section{Critical Foundations}
\label{sec:evd-foundations}
%%%%%%%%%%%%%%%%%%%%%%%%%%%%%%%%%%%%%%%%%%%%%%%%%%%%%%%%%%%%%%%%%%%%%%%%%%%%%%%

\subsection{Foundation 1: Utility Framework Requirement}

\textbf{Principle:} EBF works within an extended utility framework. Evidence must be compatible with this framework to enter CORE-WHERE.

\textbf{Implication:} Behavioral science constructs that reject utility (e.g., ``authentic self,'' ``flow state'') cannot be directly parameterized but can inform mechanism understanding.

\subsection{Foundation 2: Causal Primacy for Prediction}

\textbf{Principle:} PREDICT appendices require causal evidence. Correlations may motivate hypotheses but cannot validate predictions.

\textbf{Implication:} A PREDICT claim like ``$\theta_1 = 0.35$'' requires causal identification, not just statistical association.

\subsection{Foundation 3: Methodological Pluralism}

\textbf{Principle:} All three disciplines contribute value to EBF, but in different ways.

\textbf{Implication:} We do not rank disciplines hierarchically in general---only for specific EBF uses.


%%%%%%%%%%%%%%%%%%%%%%%%%%%%%%%%%%%%%%%%%%%%%%%%%%%%%%%%%%%%%%%%%%%%%%%%%%%%%%%
\section{Limitations and Critiques}
\label{sec:evd-critiques}
%%%%%%%%%%%%%%%%%%%%%%%%%%%%%%%%%%%%%%%%%%%%%%%%%%%%%%%%%%%%%%%%%%%%%%%%%%%%%%%

This systematization invites legitimate criticism. We document the main objections and EBF's responses:

\begin{table}[htbp]
\centering
\caption{Critiques of the BE/BS Systematization}
\label{tab:evd-critiques}
\renewcommand{\arraystretch}{1.3}
\begin{tabular}{@{}p{3.5cm}p{5cm}p{4.5cm}@{}}
\toprule
\textbf{Critique} & \textbf{Argument} & \textbf{EBF Response} \\
\midrule
\multicolumn{3}{@{}l}{\textit{\textbf{Methodological Critiques}}} \\
Economics imperialism & ``You privilege economic methods over psychological ones'' & Functional, not normative: different outputs for different purposes \\
Parameter fetishism & ``Understanding \textit{why} matters more than knowing \textit{how much}'' & Both matter; EBF needs both mechanism (BS) and magnitude (BE) \\
Laboratory paradox & ``€20 lotteries are also not real life'' & True; but real stakes $>$ hypothetical stakes for preference revelation \\
\addlinespace
\multicolumn{3}{@{}l}{\textit{\textbf{Validity Critiques}}} \\
Replication overstated & ``60--70\% still means 30--40\% don't replicate'' & Correct; individual paper assessment still required \\
WEIRD everywhere & ``Both BE and BS study Western students'' & True; external validity note applies to all Tier 1 evidence \\
Incentive crowding & ``Money distorts moral/social preferences'' & Documented (Gneezy); context-dependent discounting needed \\
\addlinespace
\multicolumn{3}{@{}l}{\textit{\textbf{Foundational Critiques}}} \\
Utility as dogma & ``Maybe utility maximization is wrong'' & EBF is pragmatic: utility works for prediction, not claimed as truth \\
Circular connectivity & ``BE is better because it connects to economics---but what if economics is wrong?'' & Fair; EBF commits to utility framework as working assumption \\
\bottomrule
\end{tabular}
\end{table}

\subsection{The Strongest Critiques}

\begin{tcolorbox}[colback=red!5!white, colframe=red!75!black,
    title=Critique 1: Parameter Fetishism]

\textbf{Argument:}
\begin{quote}
``Knowing that $\lambda = 2.25$ is less useful than understanding \textit{why} people are loss averse. Mechanisms matter more than magnitudes.''
\end{quote}

\textbf{EBF Response:}
\begin{itemize}[nosep]
    \item For \textit{understanding}: Mechanism knowledge (BS) is indeed primary
    \item For \textit{prediction}: Magnitude knowledge (BE) is necessary
    \item For \textit{intervention design}: Both are required
    \item EBF integrates both: BS for CORE-HOW, BE for CORE-WHERE
\end{itemize}

\end{tcolorbox}

\begin{tcolorbox}[colback=red!5!white, colframe=red!75!black,
    title=Critique 2: The Laboratory Paradox]

\textbf{Argument:}
\begin{quote}
``You criticize BS for hypothetical scenarios, but €20 lotteries in a Zurich lab are also not real life. Both have external validity problems.''
\end{quote}

\textbf{EBF Response:}
\begin{itemize}[nosep]
    \item Valid point---laboratory $\neq$ field for both disciplines
    \item However: incentivized choices reveal preferences better than hypotheticals
    \item Field experiments (Thaler, Gneezy) increasingly validate lab findings
    \item EBF requires external validity notes for all Tier 1 evidence
\end{itemize}

\end{tcolorbox}

\begin{tcolorbox}[colback=red!5!white, colframe=red!75!black,
    title=Critique 3: Utility Framework Commitment]

\textbf{Argument:}
\begin{quote}
``You say BE is better because it connects to welfare economics. But what if utility maximization is fundamentally wrong? Then your whole framework collapses.''
\end{quote}

\textbf{EBF Response:}
\begin{itemize}[nosep]
    \item EBF does not claim utility is \textit{true}---only that it \textit{works} for prediction
    \item Pragmatic commitment: utility framework enables quantitative policy analysis
    \item If better framework emerges, EBF structure can accommodate (replace CORE-WHAT)
    \item 50+ years of successful prediction justify working assumption
\end{itemize}

\end{tcolorbox}

\subsection{What Critics Are Right About}

\begin{enumerate}
\item \textbf{External validity is a shared problem:} Both BE and BS face generalization challenges from lab to field.

\item \textbf{Mechanisms matter:} Knowing \textit{why} enables better intervention design than knowing only \textit{how much}.

\item \textbf{Context dependence:} Parameters like $\lambda$ may vary across cultures, stakes, and domains.

\item \textbf{Replication is imperfect:} Even Tier 1 evidence should be treated with appropriate uncertainty.
\end{enumerate}

\subsection{What the Systematization Gets Right}

\begin{enumerate}
\item \textbf{Prevents false precision:} Without the hierarchy, qualitative findings get treated as calibrated parameters.

\item \textbf{Enables prediction:} Policy analysis requires numbers, not just mechanisms.

\item \textbf{Clarifies complementarity:} BE, BS, and NE each contribute---but to different questions.

\item \textbf{Protects methodological integrity:} EBF predictions are only as good as their evidence base.
\end{enumerate}

\subsection{The Balanced Position}

\begin{tcolorbox}[colback=purple!5!white, colframe=purple!75!black,
    title=EBF's Position: Functional Pluralism]

\textbf{The systematization is NOT:}
\begin{itemize}[nosep]
    \item A claim that BE is ``better'' than BS
    \item A rejection of qualitative research
    \item Economics imperialism
\end{itemize}

\textbf{The systematization IS:}
\begin{itemize}[nosep]
    \item A functional classification: different evidence for different purposes
    \item A protection against false precision
    \item A recognition that prediction requires calibrated parameters
\end{itemize}

\textbf{Key principle:}
\begin{quote}
\textit{BE asks ``how much?'' BS asks ``why?'' NE asks ``how?'' EBF needs all three---for different purposes. The systematization clarifies which evidence answers which question.}
\end{quote}

\end{tcolorbox}


%%%%%%%%%%%%%%%%%%%%%%%%%%%%%%%%%%%%%%%%%%%%%%%%%%%%%%%%%%%%%%%%%%%%%%%%%%%%%%%
\section{Computational Necessity: Why LLMs Require Evidence Classification}
\label{sec:evd-llm}
%%%%%%%%%%%%%%%%%%%%%%%%%%%%%%%%%%%%%%%%%%%%%%%%%%%%%%%%%%%%%%%%%%%%%%%%%%%%%%%

The evidence classification system is not merely an academic exercise---it is \textbf{computationally necessary} for LLM-based tools and code interpreters that power the EBF toolchain.

\subsection{The Core Problem: LLMs Have No Inherent Quality Judgment}

\begin{tcolorbox}[colback=red!5!white, colframe=red!75!black,
    title=Why LLMs Need Explicit Evidence Classification]

\textbf{LLM Properties That Create Risk:}
\begin{itemize}[nosep]
    \item \textbf{Equal treatment of all text:} An LLM treats ``Grit predicts success'' identically to ``$\lambda = 2.25$ (K\&T 1992)''
    \item \textbf{Plausible number generation:} LLMs can generate convincing but unfounded parameters
    \item \textbf{Coherence optimization:} Qualitative statements get ``translated'' to pseudo-quantitative claims
    \item \textbf{No methodological critique:} Correlation is treated like causation without explicit tagging
\end{itemize}

\textbf{The False Precision Problem at Scale:}
\begin{quote}
Without evidence classification, LLMs \textit{automate} false precision---converting qualitative insights to quantitative parameters across thousands of queries.
\end{quote}

\end{tcolorbox}

\subsection{Code Interpreters Require Numerical Input}

\begin{table}[htbp]
\centering
\caption{Why Code Interpreters Need Evidence Tiers}
\label{tab:evd-code-interpreter}
\renewcommand{\arraystretch}{1.3}
\begin{tabular}{@{}p{3.5cm}p{5cm}p{4.5cm}@{}}
\toprule
\textbf{Computational Task} & \textbf{Without Tier Classification} & \textbf{With Tier Classification} \\
\midrule
Utility calculation & $U(x) = -\lambda \cdot |x|$ fails with ``losses feel worse'' & Use $\lambda = 2.25$ (Tier 1) \\
Monte Carlo simulation & No distribution for qualitative effects & $\lambda \sim N(2.25, 0.3)$ from Tier 1 \\
Policy optimization & Cannot optimize ``fairness matters'' & Use $\alpha = 0.5$ (Fehr-Schmidt) \\
Counterfactual prediction & ``First impressions matter'' $\to$ ??? & $\theta_1 = 0.35$ (IV-identified) \\
\bottomrule
\end{tabular}
\end{table}

\subsection{The Five Critical LLM Use Cases}

\begin{enumerate}
\item \textbf{Retrieval-Augmented Generation (RAG) on Bibliography:}
\begin{itemize}[nosep]
    \item Without tiers: All 2000+ references treated equally
    \item With tiers: Filter retrieval by evidence quality (``use only Tier 1 for parameters'')
\end{itemize}

\item \textbf{LLMMC Parameter Estimation (Appendix UNMAPPED_LLM):}
\begin{itemize}[nosep]
    \item Without tiers: LLM may hallucinate parameters from qualitative sources
    \item With tiers: Restrict LLMMC input to Tier 1--2 evidence only
\end{itemize}

\item \textbf{Intervention Design:}
\begin{itemize}[nosep]
    \item Without tiers: Mechanism and magnitude evidence conflated
    \item With tiers: BS for ``why'' (Tier 3), BE for ``how much'' (Tier 1)
\end{itemize}

\item \textbf{Policy Simulation:}
\begin{itemize}[nosep]
    \item Without tiers: Pseudo-precise predictions from weak evidence
    \item With tiers: Uncertainty bands derived from evidence quality
\end{itemize}

\item \textbf{Automated Report Generation:}
\begin{itemize}[nosep]
    \item Without tiers: ``Research shows...'' without source quality
    \item With tiers: Clear attribution (``Tier 1 evidence from K\&T 1992'')
\end{itemize}
\end{enumerate}

\subsection{Practical Example: LLMMC Without vs. With Classification}

\begin{tcolorbox}[colback=yellow!5!white, colframe=yellow!75!black,
    title=LLM Parameter Extraction: Dangerous vs. Safe]

\textbf{Dangerous (No Classification):}
\begin{verbatim}
Prompt: "Estimate the effect of first impressions"
LLM: "First impressions are powerful. Studies show
      they matter a lot. I estimate ~0.4."
      ^ Where does 0.4 come from? Hallucinated.
\end{verbatim}

\textbf{Safe (With Classification):}
\begin{verbatim}
Prompt: "What Tier-1 evidence exists for
         first-exposure effect?"
LLM: "Jones et al. (2025) via IV: θ₁ = 0.32-0.36
      Ackerberg (2003) Tier 2: effect exists
      Folk wisdom: not EBF-eligible"
      ^ Clear provenance, no hallucination.
\end{verbatim}

\end{tcolorbox}

\subsection{Implementation: Evidence Tier as Metadata}

\begin{tcolorbox}[colback=blue!5!white, colframe=blue!75!black,
    title=Proposed Bibliography Tagging Schema]

Each entry in \texttt{bcm\_master.bib} should include:

\begin{verbatim}
@article{kahneman1992,
  ...
  evidence_tier = {1},
  use_for = {CORE-WHERE, parameter},
  parameter = {lambda = 2.25},
  identification = {experimental, within-subject},
  external_validity = {WEIRD, replicated}
}
\end{verbatim}

\textbf{LLM prompts can then filter:}
\begin{itemize}[nosep]
    \item ``Use only \texttt{evidence\_tier} $\leq$ 2 for parameter extraction''
    \item ``Use \texttt{evidence\_tier} = 3 for mechanism discussion only''
    \item ``Flag all \texttt{external\_validity = WEIRD} parameters''
\end{itemize}

\end{tcolorbox}

\subsection{Why This Matters for EBF's Future}

\begin{enumerate}
\item \textbf{Scalability:} Manual evidence assessment doesn't scale to 2000+ references. Tier tagging enables automated quality control.

\item \textbf{Reproducibility:} When LLMs use EBF, tier classification ensures consistent evidence use across queries.

\item \textbf{Auditability:} Tier tags create audit trails---every parameter can be traced to its evidence quality.

\item \textbf{Error Prevention:} The systematization is a \textit{guardrail} that prevents LLMs from committing the false precision error at scale.
\end{enumerate}

\textbf{Key insight:}
\begin{quote}
\textit{The evidence classification system is not optional for computational applications. Without it, LLMs will automate exactly the methodological errors this appendix warns against. The tier system is infrastructure for responsible AI-assisted behavioral economics.}
\end{quote}


%%%%%%%%%%%%%%%%%%%%%%%%%%%%%%%%%%%%%%%%%%%%%%%%%%%%%%%%%%%%%%%%%%%%%%%%%%%%%%%
\section{Open Issues}
\label{sec:evd-open}
%%%%%%%%%%%%%%%%%%%%%%%%%%%%%%%%%%%%%%%%%%%%%%%%%%%%%%%%%%%%%%%%%%%%%%%%%%%%%%%

\begin{table}[htbp]
\centering
\caption{Open Methodological Questions}
\label{tab:evd-open}
\renewcommand{\arraystretch}{1.3}
\begin{tabular}{@{}p{4cm}p{5cm}p{3cm}@{}}
\toprule
\textbf{Question} & \textbf{Why Important} & \textbf{Approach} \\
\midrule
Meta-analysis integration & How to combine multiple Tier 1 estimates? & Bayesian updating \\
Cross-cultural validity & Do BE parameters transfer to DACH? & Local replication \\
Decay over time & Do parameters remain stable? & Longitudinal tracking \\
Construct validity & When do BS constructs map to utility? & Formal translation \\
\bottomrule
\end{tabular}
\end{table}


%%%%%%%%%%%%%%%%%%%%%%%%%%%%%%%%%%%%%%%%%%%%%%%%%%%%%%%%%%%%%%%%%%%%%%%%%%%%%%%
\section{Summary}
\label{sec:evd-summary}
%%%%%%%%%%%%%%%%%%%%%%%%%%%%%%%%%%%%%%%%%%%%%%%%%%%%%%%%%%%%%%%%%%%%%%%%%%%%%%%

\begin{tcolorbox}[colback=green!10!white, colframe=green!75!black,
    title=\textbf{Appendix XX: Summary}]

\textbf{Core Question:} What evidence can EBF use, and for what purposes?

\textbf{Central Position:} This appendix classifies evidence types by function---it does \textbf{not} rank disciplines by value. Behavioral science provides essential contributions (mechanisms, hypotheses, theoretical innovation) that behavioral economics cannot provide. Both are necessary for EBF.

\textbf{Main Contribution:}
\begin{itemize}[nosep]
    \item Functional classification of BE, BS, and NE
    \item Five-tier evidence system based on methodological criteria
    \item European Laboratory Standard as quality benchmark
    \item Fehr Methodological Foundation as paradigm case
    \item Protection against false precision
    \item \textbf{Computational necessity:} Why LLMs and code interpreters require explicit evidence classification
\end{itemize}

\textbf{What Each Discipline Contributes to EBF:}
\begin{itemize}[nosep]
    \item \textbf{BE:} Calibrated parameters for CORE-WHERE ($\lambda$, $\beta$, $\theta$)
    \item \textbf{BS:} Mechanistic understanding for intervention design (``why'' questions)
    \item \textbf{NE:} Process validation for dual-system architecture (``how'' questions)
\end{itemize}

\textbf{The Fehr Recipe:} Psychological concept $\to$ Utility formalization $\to$ Incentive-compatible game $\to$ Structural estimation $\to$ Replication. This transformation (not rejection) of BS insights is the BE paradigm.

\textbf{Key Insight:}
\begin{quote}
\textit{EBF builds on 50+ years of behavioral economics research---not because economics is ``better'' than psychology, but because calibrated parameters are necessary for quantitative prediction. Behavioral science provides the ``why''; behavioral economics provides the ``how much.'' Both are essential. This appendix clarifies which evidence serves which function.}
\end{quote}

\textbf{Integration:} Protects methodological integrity while recognizing the complementary value of all three disciplines.

\end{tcolorbox}


%%%%%%%%%%%%%%%%%%%%%%%%%%%%%%%%%%%%%%%%%%%%%%%%%%%%%%%%%%%%%%%%%%%%%%%%%%%%%%%
\section{References}
\label{sec:evd-references}
%%%%%%%%%%%%%%%%%%%%%%%%%%%%%%%%%%%%%%%%%%%%%%%%%%%%%%%%%%%%%%%%%%%%%%%%%%%%%%%

All references are maintained in \texttt{bcm\_master.bib}.

\textbf{Key Methodological References:}
\begin{itemize}[nosep]
    \item Camerer, C. F. et al. (2016). ``Evaluating Replicability of Laboratory Experiments in Economics.'' \textit{Science}.
    \item Open Science Collaboration (2015). ``Estimating the Reproducibility of Psychological Science.'' \textit{Science}.
    \item Henrich, J. et al. (2010). ``The Weirdest People in the World?'' \textit{BBS}.
\end{itemize}


%%%%%%%%%%%%%%%%%%%%%%%%%%%%%%%%%%%%%%%%%%%%%%%%%%%%%%%%%%%%%%%%%%%%%%%%%%%%%%%
\section{Cross-References}
\label{sec:evd-crossref}
%%%%%%%%%%%%%%%%%%%%%%%%%%%%%%%%%%%%%%%%%%%%%%%%%%%%%%%%%%%%%%%%%%%%%%%%%%%%%%%

\noindent\textit{Related Appendices:}
Appendix HF (PREDICT-HABIT),
Appendix XIX (LIT-LEARNING),
Appendix K (LIT-FEHR),
Appendix N (LIT-KAHNEMAN),
Appendix T (LIT-CAMERER),
Appendix UNMAPPED_BBB (CORE-WHERE),
Appendix UNMAPPED_LLM (METHOD-LLMMC),
Appendix G (Glossary).

\noindent\textit{Related Chapters:}
Chapter 4 (Methodology),
Chapter 5 (Awareness),
Chapter 18 (Meta-Theory).

\noindent\textit{Master Bibliography:} \texttt{bcm\_master.bib}

\nocite{bcm_master}

% =============================================================================
% END OF APPENDIX XX
% =============================================================================

\end{document}
